\cotefile{34}
\scanpage{84}
\mspage{1}
% archive sheet number: 83
% top heading
\uline{Exposé} III \uline{le Théorème de \unsure{monodromie} par voie géométrique}

1. \uline{Comparaison de la cohomologie des \unsure{deux|2} fibres} \deleted{\unsure{spéciales}}\inserted{\uline{géométriques}} \deleted{\unsure{et|de} la fibre} \uline{spéciale et générale} : la \uline{suite exacte des cycles évanescents}.

% upper-left diagram; the manuscript strokes have no clear arrowheads
\diagram{\begin{tikzcd}[ampersand replacement=\&,column sep=small,row sep=small]
X \arrow[d,no head] \& X_K \arrow[d,no head] \& X_{\bar K} \arrow[d,no head] \\
S \arrow[dr,no head] \& K \arrow[dr,no head] \& \bar K \\
\& \widetilde S \arrow[r,no head] \& \widetilde K
\end{tikzcd}}

Notations de I 1.1. On suppose donné $X$ \uline{propre} sur $S$, d'où $X_{\unsure{\eta|y}}=X_K$ sur $\eta=\operatorname{Spec}K$, et on se propose de \deleted{\unsure{définir|déterminer}} donner des \unsure{relations} entre $H^i(X_{\bar s})$ et $H^i(X_{\bar\eta})$ (coefficients $\Lambda_\nu=\mathbb Z/\ell^{\nu+1}\mathbb Z$ ou $\mathbb Z_\ell=\varprojlim_\nu\Lambda_\nu$). On verra que cette question est liée intimement à celle des opérations de $\unsure{\pi|\Pi}$ sur $H^i(X_{\bar\eta})$ ; [par exemple on \unsure{désirerait|disposerait} un \unsure{homomorphisme|isomorphisme} $H^i(X_{\bar s})\simeq H^i(X_{\bar\eta})^\pi$].

Quitte à remplacer $S$ par $\widetilde S$, on peut supposer pour simplifier les notations que $S$ est \uline{strictement} \deleted{\unsure{hensélien}} local (et \unsure{$\simeq$|même} complet \unsure{si on y tient} \unsure{\dots}, quand $\ell\ne\operatorname{car}K$). On a
\[
\bar K=\varinjlim_{\unsure{i}}\deleted{\unsure{K}}K_\alpha\deleted{\unsure{K_\alpha}}
\tag*{(1.1)}
\]
$K_\alpha$ les sous-extensions finies de $K$, d'où
\[
X_{\bar\eta}=\varprojlim_{\unsure{i}}X_\alpha\deleted{K_\alpha}\quad\text{où}\quad X_\alpha=X_{K_\alpha},\quad\text{et}
\tag*{(1.2)}
\]
\[
H^{\unsure{*|d}}(X_{\bar\eta},\Lambda_\nu)=\varinjlim_\alpha H^*(X_{\bar K_\alpha},\Lambda_\nu)
\tag*{(1.3)}
\]
Soit $S_\alpha\inserted{=\operatorname{Spec}V_\alpha}$ le normalisé de $S$ dans $K_\alpha$, qui est encore \unsure{str.\textsuperscript{t}} local, \deleted{\unsure{soit}} et soit $s_\alpha$ le pt fermé de $S_\alpha$, $\eta_\alpha$ son pt \unsure{générique}. \inserted{\unsure{Notons}} et
\[
X_\alpha=X\times_S S_\alpha,\quad X_{\alpha0}=X_{\alpha s_\alpha}=X_0\otimes_k k_\alpha
\tag*{(1.4)}
\]
\scanpage{85}
\mspage{2}
% archive sheet number: 84
\deleted{La fibre spéciale de $X_\alpha$ est}
\[
\deleted{X_{\alpha s_\alpha}=X_0\otimes_k k_\alpha}
\]
où $X_0=X_s$, \inserted{et} $k_\alpha=k(s_\alpha)$ est une extension \inserted{finie donc} radicielle de $k$. Donc $X_{\alpha s_\alpha}\longrightarrow X_s$ induit un isomorphisme
\[
H^{\unsure{*|\alpha}}(X_s)\deleted{\Lambda_\nu)}\xrightarrow{\sim}H^*(X_{\alpha s_\alpha})\deleted{\Lambda_\nu)}
\tag*{(1.5)}
\]
% diagram in the left margin; no arrowheads visible
\diagram{\begin{tikzcd}[ampersand replacement=\&,column sep=small,row sep=small]
X \arrow[r,no head] \arrow[d,no head] \& X_\alpha \arrow[r,no head] \arrow[d,no head] \& X_\beta \arrow[d,no head] \\
S \arrow[r,no head] \& S_\alpha \arrow[r,no head] \& S_\beta
\end{tikzcd}}

D'autre part, \deleted{\unsure{les i et|les i de}} dans le diagramme commutatif
\diagram{\begin{tikzcd}[ampersand replacement=\&]
H^*(X) \arrow[r] \arrow[d,"\sim"] \& H^*(X_\alpha) \arrow[d,"\sim"] \\
H^*(X_s) \arrow[r] \& H^*(X_{\alpha,s_\alpha})
\end{tikzcd}\quad (1.6)}
les flèches verticales sont \unsure{également} des isomorphismes, car $X/S$ \unsure{donc|et} $X_\alpha/S_\alpha$ \uuline{\unsure{propre|propres}} sur $S$. Donc
\[
H^*(X)\xrightarrow{\sim}H^*(X_\alpha)
\tag*{(1.7)}
\]

Considérons maintenant la suite exacte de cohomologie relative sur $X_\alpha$, muni de l'ouvert $X_{{\unsure{\eta|y}}_\alpha}\deleted{\unsure{K_\alpha}}$
\[
\cdots\longrightarrow H^i_{X_{\alpha0}}(X_\alpha,\Lambda_\nu)\longrightarrow H^i(X_\alpha,\Lambda_\nu)\longrightarrow H^i(X_{\eta_\alpha},\Lambda_\nu)\longrightarrow H^{i+1}_{X_{\alpha,0}}(X_{\deleted{\unsure{\eta}}\alpha},\Lambda_\nu)\longrightarrow\cdots
\tag*{(1.8)}
\]
Passant à la limite sur $\alpha$, et utilisant (1.3) et (1.7), on trouve une suite exacte
\[
\cdots\longrightarrow{\unsure{\Phi|\varphi}}^i_\nu\longrightarrow H^i(X_s,\Lambda_\nu)\longrightarrow H^i(X_{\bar\eta},\Lambda_\nu)\longrightarrow\Phi^{i+1}_\nu\longrightarrow\cdots
\tag*{(1.9)}
\]
où on a
\[
\Phi^i_\nu=\varinjlim_\alpha H^i_{X_{\alpha0}}(X_\alpha,\Lambda_\nu)
\tag*{(1.10)}
\]
\scanpage{86}
\mspage{3}
% archive sheet number: 85
1.11. Le groupe ${\unsure{\Phi|\varphi}}^i_\nu$ peut être appelé \uline{groupe des cycles évanescents en dimension $i$} (pour $X/S$\inserted{et les coeff.\ $\Lambda_\nu$}) et la suite exacte (1.9) la \uline{suite exacte des cycles évanescents}. Les cycles évanescents apparaissent ici comme mesurant ``l'écart'' entre la cohomologie de $X_s$ et de $X_{\bar{\unsure{\eta|y}}}$, fibres géométriques spéciale et générale.

Passant à la limite \inserted{projective} sur $\nu$, on trouve (les groupes de la suite (1.9) étant \uline{finis} \unsure{$\ell$)|(!)} par le th. de finitude)
\[
\cdots\longrightarrow\Phi^i\longrightarrow H^i(X_s,\mathbb Z_\ell)\longrightarrow H^i(X_{\bar\eta},\mathbb Z_\ell)\longrightarrow\Phi^{i+1}\longrightarrow\cdots,
\tag*{(1.\deleted{\unsure{9}}\deleted{\unsure{9}}9 bis)}
\]
avec
\[
\Phi^i=\varprojlim_\nu\Phi^i_\nu=\varprojlim_\nu\varinjlim_\alpha H^i_{X_{\alpha0}}(X_\alpha,\Lambda_\nu).
\tag*{(1.10 bis)}
\]

\deleted{1.12.} On appelle aussi (1.9 bis) suite exacte des cycles évanescents pour les coeff.\ $\mathbb Z_\ell$, et $\Phi^i$ le \unsure{groupe} des cycles évanescents en dim.\ $i$ pour les coeff.\ $\mathbb Z_\ell$.

1.12. \uline{Passage du local au global} Pour calculer les $\Phi^i_\nu$, \deleted{il y a lieu} utilisant la \deleted{\unsure{repr}}\unsure{expression} (1.10), on regarde pour chaque $\alpha$ la suite spectrale
\[
H^*_{X_{\alpha0}}(X_\alpha,\Lambda_\nu)\Longleftarrow E_2^{pq}=H^p(X_{\alpha0},\underline H^q_{X_{\alpha0}}(\Lambda_{\nu X_\alpha}))
\tag*{(1.13)}
\]
Ces suites spectrales, pour $\alpha$ variable, forment un système inductif, et donnent à la limite, compte tenu que $X_{\alpha0}\longrightarrow X_0$ induit une équivalence de \unsure{sites étales}
\scanpage{87}
\mspage{4}
% archive sheet number: 86
\[
{\unsure{\Phi|\varphi}}^*_\nu\Longleftarrow E_2^{pq}=H^p(X_0,\underline\Phi^q_\nu)
\tag*{(1.14)}
\]
où $\underline\Phi^q_\nu$ est un \uline{faisceau} sur $X_0$ qui, via identification de $X_{\alpha0}$ avec $X_0$, peut s'écrire
\[
\underline\Phi^q_\nu=\varinjlim_\alpha\underline H^*_{X_{\alpha0}}(\Lambda_{\nu X_\alpha})
\tag*{(1.15)}
\]
Le calcul de ces faisceaux sur $X_0$ est \uline{de nature locale sur} $X$. [Sous réserve \deleted{que} de vérifications de passages à la limite, on conclut de \unsure{même}
\[
\Phi^*\Longleftarrow E_2^{pq}=H^p(X_0,\underline\Phi^q)
\tag*{(1.14 bis)}
\]
avec
\[
\underline\Phi^q=\varprojlim_\nu\underline\Phi^q_\nu,
\tag*{(1.15 bis)}
\]
\inserted{mais} que nous n'utiliserons pas tels quels.\deleted{]}\inserted{Il faudrait en effet \unsure{supposer}, \unsure{voire justifier ce point de vue}, que les $\underline\Phi^q_\nu$ sont \unsure{constructibles|constructible} et que le syst. projectif $(\underline\Phi^q_\nu)_{\nu\geq1}$ définit un faisceau $\ell$-adique constructible sur $X$\unsure{\dots}).}

1.16. On trouve ainsi un principe de localisation fort utile pour comparer la cohomologie de $X_s$ et de $X_{\bar{\unsure{\eta|y}}}$. Notons par exemple, \deleted{\unsure{supp}}\deleted{que le théo} (Supposant que $\ell\ne p$) que \deleted{\unsure{avec}}\inserted{si $X\longrightarrow S$ est \uline{lisse}} le théorème de pureté pour l'inclusion $s\longrightarrow S$, et le th. de changement de base lisse, nous permettent de calculer $\underline H^{\deleted{\unsure{*}}i}_{X_{\alpha0}}(\Lambda_{\nu X_\alpha})$ aux points où $X_\alpha$ est \uline{lisse} sur $S_\alpha$ : c'est nul pour $i\ne0$, et \deleted{\unsure{\simeq}} $\simeq(\mu_{\ell^\nu}^{\otimes(-1)}{}_{X_{\alpha0}})$ \deleted{\unsure{autrement}}\inserted{si $i=2$}.

On en conclut que $\underline\Phi^i_\nu=0$ si $i\ne2$, et $\underline\Phi^2_\nu=\deleted{\unsure{H}}\varinjlim_\alpha(\mu_{\ell^\nu})_{X_{\alpha0}}$. Il faut déterminer les
\scanpage{88}
\mspage{5}
% archive sheet number: 87
morphismes de transition, \inserted{entre ces faisceaux, \unsure{identifiés comme} le faisceau $\mu_{\ell^\nu}$ sur $X_0$} On trouve que si $S_{\deleted{\unsure{\beta}}\alpha}$ est de degré $n_\alpha$ sur $S$, donc $S_\beta$ de degré $\inserted{n_{\alpha\beta}=}n_\beta/n_\alpha$ sur $S_\alpha$ (si $\beta\geq\alpha$), alors le morphisme de transition est la multiplication par $n_{\alpha\beta}$. \inserted{(On est ramené à faire \unsure{cette} vérification pour $X=S$\dots)} On en conclut aussitôt que $\underline{\unsure{\Phi|\varphi}}^2_\nu=\varinjlim(\mu_{\ell^\nu})$ \deleted{\unsure{est isomorphe}}\inserted{est également nul}. \deleted{\unsure{plus simplement}} \uline{Donc} $\underline\Phi^i_\nu=0$ \uline{pour \unsure{tt|$\forall$} $i,\nu$}, \deleted{d'où} \deleted{grâce au fait} On en déduit grâce à (1.14) que $\Phi^i_\nu=0$ pour tt $i,\nu$, puis, grâce à (1.9) que
\[
H^i(X_s,\Lambda_\nu)\xrightarrow{\sim}H^i(X_{\bar{\unsure{\eta|y}}},\Lambda_\nu)
\tag*{(1.16.1)}
\]
($X$ lisse sur $S$, \deleted{\unsure{supposant}} $\ell\ne p$) : \deleted{c'est}\inserted{on \deleted{\illegible{1}} \unsure{retrouve}} le \uline{th. de spécialisation} de la cohomologie pour un morphisme propre et lisse. (SGA 4 $\overline{\underline{\mathrm{XVI}}}$ \dots). [Bien entendu, la \unsure{même} démonstration marcherait avec des coefficients constants tordus).

Compte tenu de la nature locale des $\underline\Phi^i_\nu$ en général, on conclut de ce qui précède la

\uline{Proposition} 1.17. Soit $T$ l'ensemble des $x\in X_0$ en lesquels $f\inserted{:X\longrightarrow S}$ n'est pas lisse, \inserted{et supposons $\ell\ne p$.} Alors les faisceaux $\underline\Phi^i_\nu$ des cycles évanescents sont à support $\subset T$.

\deleted{les points de $T$ étant \unsure{fermés},}\deleted{\unsure{sur}}

1.18. Supposons p.ex. que $T$ soit \uline{discret}. Alors dans \deleted{$\unsure{E_2^{pq}}$} (1.14) on a $E_2^{pq}=0$ pour $p\ne0$, donc on trouve des isomorphismes
\scanpage{89}
\mspage{6}
% archive sheet number: 88
\[
{\unsure{\Phi|\varphi}}^i_\nu=\sum_{t\in T}(\underline\Phi^i_\nu)_t,\quad\deleted{\text{avec}}
\tag*{(1.18)}
\]
où
\[
(\underline\Phi^i_\nu)_t=\varinjlim_\alpha\underline H^i_{X_{\alpha0}}(\Lambda_{\nu X_\alpha})_{t_\alpha}
\tag*{(1.19)}
\]
où $t_\alpha$ est l'unique point de $X_\alpha$ au dessus de $t$.

On voit donc que dans le cas où l'ensemble $T$ des pts de $X_0$ en lesquels $f$ est non lisse est \uline{fini}, le \deleted{calcul \unsure{du|des}} groupe $\Phi^i_\nu$ des cycles évanescents est \inserted{une somme de termes} \uline{purement locaux} \deleted{\unsure{en les points de}}\inserted{(au sens de la top. étale)} relatifs aux différents points de $\unsure{t|T}$). \deleted{\unsure{1.19}} Nous donnerons plus tard des exemples \deleted{de calculs} d'application de ce principe à des cas particuliers intéressants.

\inserted{1.20.} Mais nous allons nous tourner \inserted{d'abord} vers une \deleted{autre}\inserted{variante} \deleted{méthode} de la méthode précédente pour \deleted{façon de} formuler les relations entre la cohomologie de \deleted{$\unsure{\mathcal E|E}$} $X_s$ et de $X_{\bar{\unsure{\eta|y}}}$, \inserted{qui sera} plus commode dans la démonstration du th. de monodromie.

Notons toutefois \inserted{(\unsure{tout de suite})} que \inserted{On constate aisément que \unsure{.|:|;}} par construction même, le \deleted{\unsure{2.}} groupe d'inertie $I=\operatorname{Gal}(\bar K/K)$ opère sur les groupes $\Phi^i_\nu$ et \deleted{\unsure{les}} sur les faisceaux $\underline\Phi^i_\nu$, que la suite spectrale (1.14) peut être interprétée comme une suite spectrale de $I$-modules, et la suite exacte (1.9) comme une suite exacte de $I$-modules, étant entendu que $I$ opère trivialement sur $H^*(X_s,\Lambda_\nu)$. \inserted{De \unsure{même} après passage à la limite.} Il sera possible de \unsure{prouver} le th. de
\scanpage{90}
\mspage{7}
% archive sheet number: 89
monodromie, dans les cas \deleted{\unsure{qui seront}} traités plus bas, en prouvant que l'action de $I$ sur chacun des ${\unsure{\Phi|\varphi}}^i$ est \deleted{quasi}\margin{\unsure{non!}} unipotente, \deleted{et}\inserted{en prouvant} pour ceci, que \deleted{suffira de} l'action sur chacun des \uline{faisceaux} $\ell$-adiques $\underline\Phi^i$ est \deleted{\unsure{q}} \uline{essentiellement triviale}, et mieux, qu'il existe un sous-groupe \deleted{\unsure{d'ind}} ouvert fixe $I'$ de $I$ qui opère trivialement sur \uline{tous} les \deleted{$\unsure{\Phi}$} $\underline\Phi^i_\nu$ ($i,\nu$ variables).

\deleted{Les calculs \unsure{précédents}}. C'est là un résultat de nature \uline{purement local} sur $X$, qui est accessible au calcul sous des conditions favorables, \unsure{examinées} au n° suivant.

1.21. Généralisation des développements précédents au cas d'un ouvert $U$ de la fibre générique. \inserted{($X$ tjs propre sur $S$). Soit $Y=X-U\supset X_0$.} On trouve une suite exacte des cycles évanescents
\[
\cdots\longrightarrow\Phi^i_\nu\longrightarrow H^i(X_s,\Lambda_\nu)\longrightarrow H^i(U,\Lambda_\nu)\longrightarrow\Phi^{i+1}_\nu\longrightarrow\cdots
\tag*{(1.22)}
\]
avec
\[
\Phi^i_\nu=\varinjlim_\alpha H^i_{Y_\alpha}(\deleted{\unsure{X}}X_\alpha,\Lambda_\nu),
\tag*{(1.23)}
\]
où $Y_\alpha$ est l'image inverse de $Y_\alpha$ dans $X_\alpha$. La suite spectrale de localisation est
\[
\Phi^*_\nu\Longleftarrow E_2^{pq}=H^p(X_0,\underline\Phi^q_\nu),
\tag*{(1.24).}
\]
\[
\underline\Phi^q_\nu=\varinjlim_\alpha\bigl(\underline H^q_{Y_\alpha}((\Lambda_\nu)_{X_\alpha})\mid X_{\alpha0}\bigr).
\tag*{(1.25)}
\]
\scanpage{91}
\mspage{8}
% archive sheet number: 90
\deleted{\unsure{1.25}}

1.26. \uline{Restriction aux revêtements modérément ramifiés} de $S$.

Comme on verra par la suite, le calcul des $\underline H^i_{\unsure{X_{\alpha0}|Y_\alpha}}(\Lambda_{\nu X_\alpha})$ peut se faire, parfois, à condition de se borner aux $S_\alpha$ qui sont modérément ramifiés sur $S$, correspondant donc aux \deleted{les} sous-extensions \inserted{finies} de l'extension $\bar K_t\subset\bar K$, \deleted{\unsure{leur}} correspondant au $p$-sous-groupe de Sylow $P$ de $I$. \deleted{\unsure{Puisque}} \deleted{Dans ce cas on voit} \deleted{La suite} La suite exacte de \unsure{Hochschild-Serre} pour $\deleted{X}U_{\bar K}\longrightarrow\deleted{X}U_{\bar K_t}$ \inserted{de gr $P$} donne alors
\[
H^{\unsure{*|\alpha}}(U_{\bar K_t},\Lambda_\nu)\Longleftarrow E_2^{p,q}=H^p(P,H^q(U_{\bar K},\Lambda_\nu)),
\]
et on a $E_2^{pq}=0$ si $p\ne0$, car $H^q(U_{\bar K},\Lambda_\nu)$ est de $\ell$-torsion, $\ell\ne p$. Donc on trouve
\[
H^*(U_{\bar K_t},\Lambda_\nu)=H^*(U_{\bar K},\Lambda_\nu)^P,
\tag*{(1.27)}
\]
invariants sous le groupe de $p$-Sylow. D'ailleurs, le premier \unsure{est aussi|est encore} isomorphe à $\varinjlim_\alpha H^*(U_{K_\alpha},\Lambda_\nu)$, où $K_\alpha$ parcourt les sous-extensions finies \uline{de} $\bar K_t$, \deleted{de} \unsure{ou|à} si on veut, les
\[
K(n)=K({\unsure{\pi|\Pi}}^{1/n})\qquad\text{pour }(p,n)=1.
\tag*{(1.28)}
\]
\scanpage{92}
\mspage{9}
% archive sheet number: 91
\margin{\unsure{\S} On trouve alors une suite exacte de la forme (1.29), mais où on peut écrire $H^i(U_{\bar{\unsure{\eta|y}}},\mathbb Z_\ell)$ au lieu de $H^i(U_{\bar\eta},\mathbb Z_\ell)^P$.}

On trouve alors une suite exacte
\[
\longrightarrow({\unsure{\Phi|\varphi}}^i_\nu)_{\mathrm{tame}}\longrightarrow H^i(X_s,\Lambda_\nu)\longrightarrow H^i(U_{\bar\eta},\Lambda_\nu)^P\longrightarrow(\Phi^{i+1}_\nu)_{\mathrm{tame}}\longrightarrow\cdots
\tag*{(1.29)}
\]
avec
\[
(\Phi^i_\nu)_{\mathrm{tame}}=\varinjlim_{(n,p)=1}H^{\unsure{0|i}}_{Y_n}(X_n,\Lambda_\nu),
\tag*{(1.30)}
\]
où $X_n=X\times_S S(n)$, $S(n)$ le normalisé de $S$ dans $K(n)$, $Y_n$ l'image inverse de $Y=X-U$ dans $X(n)$. La suite spectrale de localisation est ici
\[
(\Phi^*_\nu)_{\mathrm{tame}}\Longleftarrow E_2^{pq}=H^p(X_0,(\underline\Phi^q_\nu)_{\mathrm{tame}})
\tag*{(1.31)}
\]
avec
\[
(\underline\Phi^q_\nu)_{\mathrm{tame}}=\varinjlim_{(n,p)=1}\bigl(\deleted{\unsure{L}}\underline H^q_{Y(n)}((\Lambda_\nu)_{X(n)})\bigr)\mid X(n)_0,
\tag*{(1.32)}
\]
où $X(n)_0$ est la \deleted{le \unsure{revêtement de}} fibre spéciale de $X(n)$, identifiée (comme \deleted{\unsure{objet}}\inserted{site} étale) à $X_0\deleted{\unsure{\simeq}}$.

\inserted{Cette suite \deleted{spectrale}\inserted{exacte} (1.29)\unsure{\S} devient \uline{utile si $P$ opère trivialement} sur $H^i(U_{\bar\eta},\mathbb Z_\ell)$, \deleted{\unsure{néanmoins}} condition qui sera vérifiée en tout cas après extension finie de $S$ (cf \unsure{I}).}
\margin{et celle qui s'en déduit par $\varprojlim_\nu$}
\margin{(\unsure{l'hyp. de déf.|l'hyp. de départ} sur $\mathbb Z_\ell$)}

1.33. Les développements restent valables \inserted{mutatis mutandis} si on remplace $P$ par n'importe quel sous-groupe invariant fermé $Q$ de $I$ \deleted{tel que $I$} qui est ``premier à $\ell$'', i.e. tel que $I/Q$ couvre le quotient $\mathbb Z_\ell(1)$ de $I_t$, composante $\ell$-primaire de $I/P$. Un cas particulier \deleted{\unsure{intéressant}} qui peut être utile est celui où $I/Q\simeq\mathbb Z_\ell(1)$ est cette composante $\ell$-primaire, ce qui signifie qu'on fait parcourir à $K_\alpha$ les \deleted{\unsure{de}} $K_{\unsure{\alpha}}(n)$ avec $n$ de la forme $\ell^r$. \deleted{Dans ce cas}
\scanpage{93}
\mspage{10}
% archive sheet number: 92

2. \uline{Comparaison de la cohomologie de la fibre spéciale et générale par la suite spectrale de Leray.}

2.1. $X$ propre sur $S$, $U$ ouvert de $X_{\unsure{\eta|y}}$.

La suite spectrale de \deleted{localisation} \inserted{cohomologie relative} sera ici remplacée, essentiellement, par la suite spectrale \inserted{(légèrement différente)} de l'inclusion ouverte $U_\alpha\longrightarrow X_\alpha$, ou plutôt par la limite de ses suites spectrales. \deleted{\unsure{La plus}} Nous ne supposons pas ici que $k(s)$ soit séparablement clos, car nous désirons \deleted{\unsure{générale}} \inserted{parfois} suivre les opérations de $\unsure{\pi|\Pi}$ tout entier, pas de $I$ seulement. Cela nous amène à introduire, en plus de l'ouvert $U$ de $X_\eta$ et de $\bar U$, $\bar X=X_{\bar\eta}$, les schémas \unsure{analogues|conjugués} $\tilde X$ et $\tilde U$ sur $\tilde S$.

\diagram{\begin{tikzcd}[ampersand replacement=\&]
U \arrow[d,phantom,"\cap" description] \& \tilde U \arrow[d,phantom,"\cap" description] \& \bar U=U_{\bar\eta} \arrow[d,phantom,"\cap" description] \\
X \arrow[d,no head] \& \tilde X \arrow[l] \arrow[d,no head] \& \bar X=X_{\bar\eta} \arrow[l,"\unsure{g|\beta}"] \arrow[d,no head] \\
S \& \tilde S \arrow[l] \& \bar S \arrow[l,no head]
\end{tikzcd}\quad (2,2)}

\deleted{On utilise la suite} Le groupe $\pi$ opère sur la situation
\diagram{\begin{tikzcd}[ampersand replacement=\&]
\tilde U \arrow[d,phantom,"\cap" description] \& \bar U \arrow[l] \arrow[d,phantom,"\cap" description] \\
\tilde X \& \bar X \arrow[l]
\end{tikzcd}}
le sous-groupe $I$ opérant par $\tilde X$-automorphismes.

Nous considérons la suite spectrale de Leray pour le composé \deleted{$\bar U$}
\[
\tilde X\xleftarrow{g}\bar U\ ,
\tag*{(2,3).}
\]
on trouve
\[
H^*(\bar U,\Lambda_\nu)\Longleftarrow E_2^{pq}=H^p(\tilde X,{\unsure{R|\mathbb R}}^qg_*((\Lambda_\nu)_{\bar U})\ .
\tag*{\deleted{\unsure{(2.3)}} (2.4)}
\]
\deleted{On peut considérer cette suite spectrale comme une suite}
Le th de changement de base propre permet de l'écrire sous la forme
\scanpage{94}
\mspage{11}
% archive sheet number: 93

\[
H^{\unsure{*|\alpha}}(\bar U,\Lambda_\nu)\Longleftarrow E_2^{pq}=H^p(\tilde X_0,\underline\Psi_\nu^q)
\tag*{(2.5)}
\]
avec
\[
\underline\Psi_\nu^q=R^qg_*((\Lambda_\nu)_{\bar U}\mid\tilde X_0)=\varinjlim_\alpha R^qg_{\alpha*}((\Lambda_\nu)_{U_\alpha})\mid\tilde X_0,
\tag*{(2.6)}
\]
où on a écrit $\bar K=\varinjlim K_\alpha$ ($K_\alpha$ extensions finies de $\bar K/\tilde K$), $g_\alpha:\deleted{\unsure{U_\alpha}}\tilde U\mathbin{\times}_{\tilde S}S_\alpha\longrightarrow\tilde X$ \deleted{\unsure{et}} l'immersion canonique.

\deleted{\unsure{2.7.}} 2.7. Ici le calcul des faisceaux $\underline\Psi_\nu^q$ est de nature locale sur $\tilde X$. Notons d'ailleurs \unsure{leurs|les} relations de ces faisceaux avec \deleted{\unsure{ceux de}} les faisceaux $\underline{\unsure{\Phi|\varphi}}_\nu^q$ du n° 1 (que nous n'utilisons pas)
\[
\left\{\begin{aligned}
&\underline\Phi_\nu^q=\underline\Psi_\nu^{q-1}\qquad\text{si}\quad q\geq2\\
&0\longrightarrow\underline\Phi_\nu^0\longrightarrow(\Lambda_\nu)_{X_0}\longrightarrow\underline\Psi_\nu^0\longrightarrow\underline\Phi_\nu^1\longrightarrow0\ .
\end{aligned}\right.
\tag*{(2.8)}
\]
Donc : \unsure{peu de chose près}, ces faisceaux $\underline\Psi_\nu^{\unsure{i|q}}$ sont les faisceaux de cycles évanescents pour $\tilde X/\tilde S$ et l'ouvert $\bar U$ de $\tilde X_{\tilde{\unsure{\eta|y}}}$ (coeff les $\Lambda_\nu$).

2.9. Il y a lieu de considérer (2.5) comme une suite spectrale de $\unsure{\pi|\Pi}$-modules, et l'expression (2.6) comme un \unsure{nom|hom|isom.} de faisceaux \inserted{sur $\tilde X_0$} où $\pi$ opère, (de façon compatible avec ses opérations sur $\tilde X_0$ via $\pi_0$). Nous nous intéressons d'abord à l'action du sous-groupe \deleted{$I$} d'inertie $I$.

2.10. On peut donner encore une variante évidente des développements précédents, en remplaçant $\bar K$ par $\bar K_t$, cf 1.26. On trouve
\scanpage{95}
\mspage{11'}
% archive sheet number: 94

\[
H^{\unsure{*|\alpha|q}}(\bar U,\Lambda_\nu)^P\Longleftarrow E_2^{pq}=H^p(\tilde X_0,(\underline\Psi_\nu^q)_{\mathrm{tame}})\ ,
\tag*{(2.\unsure{5})bis)}
\]
avec
\[
\begin{aligned}
(\underline\Psi_\nu^q)_{\mathrm{tame}}&=R^qg_{t*}^{\deleted{\unsure{\mathrm{tame}}}}((\Lambda_\nu)_{U_{\bar K_t}})\mid\tilde X_0\\
&=\varinjlim_{(n,\unsure{\ell|p})=1}R^qg_{n*}(\Lambda_\nu)_{U_{K(n)}}\mid\tilde X_0\ ,
\end{aligned}
\tag*{(2.\unsure{6} bis)}
\]
où
\[
K(n)=\tilde K({\unsure{\pi|\Pi}}^{1/n})\subset\bar K_t\unsure{,|;}\subset\bar K
\]
et
\[
g_n:U_{K(n)}\longrightarrow\tilde X
\]
est la projection canonique.
\scanpage{96}
\mspage{11''}
% archive sheet number: 95

\uline{Remarque} \unsure{2.7 bis|2.7}. Il semble plausible qu'on doive pouvoir montrer, \deleted{\unsure{à}} \inserted{lorsqu'on dispose de} la résolution des singularités, que les faisceaux $\underline{\unsure{\Phi|\varphi}}_\nu^{(q)}$ \inserted{(ou les $\underline\Psi_\nu^{(q)}$, c'est pareil)}, \deleted{\unsure{sera}} sont constructibles si $\ell\ne\deleted{\unsure{\mathrm{car}}}p$, et que pour $\nu$ variable, le système projectif des $(\underline\Phi_\nu^{(q)})$ \inserted{(ou $\underline\Psi_\nu^{(q)}$)} \deleted{\unsure{est} définit de façon} admet une limite projective dans la catégorie des faisceaux $\ell$-adiques constructibles. S'il en était ainsi, on déduirait par passage à la limite \inserted{projective}, à partir des suites spectrales \deleted{\unsure{(1.14)}} \inserted{(\deleted{\unsure{III}}1.14)} resp. \deleted{\unsure{III}} (2.5), des suites spectrales \deleted{de faisceaux} \inserted{\unsure{à composantes cohomologiques|à composantes cohomologique}} $\ell$-adiques, \unsure{\S} C'est ce \unsure{qu'on|qui} \unsure{constatera|constate} dans un cas particulier typique au n° \unsure{5}, -- et il est plausible que le cas général puisse s'y ramener par les \uline{arguments} \unsure{standart|standard}, via la résolution des \uline{singularités}.
\scanpage{97}
\mspage{12}
% archive sheet number: 96

3. \uline{Cas} \deleted{où} \inserted{$X$ \uline{régulier et}} $Y=X-U$ \deleted{\unsure{est}} \uline{un diviseur à croisements normaux.} \inserted{\uline{Enoncé du th. principal et de ses corollaires}}

3.1 \inserted{On suppose $\ell\ne p$.}
\deleted{(Dans le \unsure{résultat qui suit}, \unsure{nous n'indiquons que|nous indiquons que} l'action du gr d'inertie, nous supposons \inserted{\unsure{(d. la réd. définit.)}} pour simplifier les notations que $S$ est \uline{local}.)}
(\unsure{Nous} \unsure{noterons} \deleted{cependant dans le cas général} que les hypothèses : $X$ régulier, $Y$ diviseur à croisements normaux, sont stables par \deleted{\unsure{passage de} $\unsure{X}\unsure{\tilde X}$} \inserted{changt. de base $\tilde S\mathbin{-}S$}).

\inserted{Nous supposons de plus que tout schéma régulier \unsure{quasifini} sur $X$ satisfait au th. de pureté cohomologique (cf 4.1).}
% slanted insertion from the left margin
Nous supposons \unsure{pour l'ins. seulement|pour l'instant seulement} que $H^i(U_{\bar K},\mathbb Z/\ell\mathbb Z)$ est fini, ce qui est le cas en tous cas si $U$ est \unsure{propre lisse} sur $K$ (\unsure{est|et} un schéma régulier).
\deleted{\unsure{résultat} \illegible{3} \unsure{de schéma} \illegible{3} $X$.}
\deleted{Nous supposons de plus que \illegible{5} \unsure{l'hypothèse} \illegible{2} $\Lambda_\nu$.}
\deleted{\unsure{rég.} \illegible{3}) \unsure{satisfait au th. de pureté} \illegible{1}}
\margin{\deleted{\unsure{Réd. définitive}}}
\margin{\unsure{de 1 et 3}}

\uline{Théorème} 3.2. Avec les hypothèses du titre, \deleted{et} \deleted{et, pour $i$, $H^i(U_{\bar K},\mathbb Z/\ell\mathbb Z)$ fini} \inserted{de 3.1} \deleted{\unsure{$S$ étant clos séparément}} $H^{\unsure{i|*}}(U_{\bar K},\mathbb Z_\ell)$ est un $\mathbb Z_\ell$-module de type fini et l'action du \inserted{groupe d'inertie} $I$ sur \inserted{$E=H^i(U_{\bar K},\mathbb Z_\ell)^P$} est essentiellement unipotente. \deleted{\unsure{Pour}} ($P$ étant le $p$-sous-groupe de Sylow de $I$, cf I\unsure{\ldots}).

Plus précisément, \deleted{\unsure{soit}} posons (
\[
\deleted{\unsure{X_0}}\qquad X_0=\sum m_{\unsure{\lambda|A}}X_0^{(\lambda)}
\]
\deleted{\unsure{où}}
(identité de diviseurs sur $X$), \deleted{et} où les $X_0^{(\lambda)}$ sont les composantes irréductibles de $X_0$, et les $m_\lambda>0$. Soit $m$ le ppcm des $m_\lambda$, \deleted{et} $m'$ la composante première à $p$ de $m$, et $I'$ le sous-groupe \inserted{ouvert} d'indice $m'$ de $I$. Alors l'action de $I'$ sur $E=\deleted{\unsure{H^i(U_{\bar K},\mathbb Z_\ell)}}\inserted{H^i(U_{\bar K},\mathbb Z_\ell)^P}$ est \uline{unipotente d'échelon $i+1$}, i.e. il y a une filtration \inserted{pure} de $E$ \inserted{par des $\mathrm{Fil}^p$}, invariante sous $I$, ayant $\unsure{p|i}+1$ quotients, et l'action de $I$ sur les $\unsure{\mathrm{Gr}^p|\mathrm{Gr}_F^p}(E)$ étant triviale. En d'autres termes,
\scanpage{98}
\mspage{13}
% archive sheet number: 97

si $g$ est un générateur topologique de $I_t=I/P$ (\unsure{horreur!}), \inserted{$g_E$ son action sur $E$,} on a la relation
\[
(\mathrm{id}-g_E^{m'})^{i+1}=0
\tag*{(3.3)}
\]

\uline{Remarques} 3.4. a) La démonstration que nous allons donner \inserted{au \deleted{\unsure{4}}5} \unsure{marche|montre} également, mutatis mutandis, dans le cas analytique complexe, pour un morphisme propre $f:X\mathbin{-}S$ d'espaces analytiques, $S$ régulier de dim 1. \inserted{(auquel cas il n'y a plus lieu d'introduire \unsure{le groupe} $P$!)} L'énoncé dans ce cas est dû à P.A. Griffiths.
\margin{plus haut}
% slanted; bracket alongside paragraph a)

b) La démonstration qui suit \deleted{\unsure{montre que}} \inserted{donne en} fait une filtration \uline{canonique} de $E$, qui sera de plus invariante sous l'action de $\unsure{\pi|\Pi}$ tout entier, pas seulement de $I$. Elle sera étudiée plus bas à un point de vue arithmétique.

Signalons tout de suite :

\uline{Corollaire 3.5} \deleted{Supposons que} Soit \deleted{$\unsure{X}$} \inserted{$n\geq0$, et} \deleted{$U$ lisse quasi-projectif de type fini sur $K$\inserted{de dim \deleted{relative} $n$},} supposons \deleted{que}

(i) Pour tte extension finie $K'$ de $K$, si $S'$ est le normalisé de $S$ dans $K'$, la \uline{résolution des singularités forte} est possible pour \unsure{tt|$\forall$} schéma de type fini $X'$ sur $\hat S'$ de dim relative $\leq n$.

(ii) Pour tout \deleted{$\unsure{X'|K'},\unsure{S'}$} $X'$ comme dessus, avec $X'$ régulier, et tout diviseur régulier $Y'$ de $X'$, le th.
\scanpage{99}
\mspage{14}
% archive sheet number: 98

de pureté cohomologique est vrai pour le couple $(X',Y')$ et pour $\ell$.

Sous ces conditions, \inserted{soit $U$ lisse quasi-projectif sur $K$, et} pour \unsure{tt|$\forall$} entier $i$, l'action du groupe d'inertie $I$ sur $\deleted{\unsure{H^i(U_{\bar K})}}H^i(Y_{\bar K},\mathbb Z_\ell)$ est \deleted{\unsure{quasi}}\inserted{\unsure{essentiellement}} unipotente d'échelon $i+1$.

\uline{Dém} En vertu de I\ldots, on peut supposer $S=\hat S$. \deleted{On peut que} Comme $U$ est quasi-projectif sur $K$, \inserted{de dim relative $n$,} \deleted{il provient d'un} $\exists X$ \deleted{projectif} propre sur $S$ (et = projectif sur $S$), \deleted{plat sur $S$} de dim relative $n$. On peut supposer $X=$ adhérence schématique de $U$, donc $X$ plat sur $S$ et $X$ réduit. \deleted{Notons} \inserted{Rappelons} que l'hypothèse de \unsure{lissité} \deleted{de} sur $U$ implique $H^i(U_{\bar K},\mathbb Z/\ell\mathbb Z)$ fini, donc $E=H^i(U_{\bar K},\mathbb Z_\ell)$ de t.f. sur $\mathbb Z_\ell$.

Quitte : faire une extension finie \inserted{séparable} sur $K$, on peut supposer que l'action de $P$ sur $E$ est triviale (cf I\ldots). Utilisant la résolution \deleted{des} forte des singularités \inserted{(hyp. (i))}, on est ramené au cas $X$ régulier, et $Y=X-U$ un diviseur à croisements normaux. On peut alors appliquer 3.2 grâce à l'hypothèse (ii), et on gagne.

\uline{Remarque} 3.6. Si on admet le résultat de compactification de Nagata, la démonstration précédente marche
